\documentclass[11pt,a4paper]{article}

%========================
% PACKAGES
%========================
\usepackage[utf8]{inputenc}
\usepackage[T1]{fontenc}
\usepackage[french]{babel}
\usepackage[a4paper,margin=1.8cm]{geometry}
\usepackage{amsmath,amssymb}
\usepackage{lmodern}
\usepackage{fancyhdr}
\usepackage{enumitem}

\pagestyle{fancy}
\fancyhf{}
\fancyhead[L]{3\ieme}
\fancyhead[C]{TD – Mécanique}
\fancyhead[R]{\thepage}

\title{\Huge TD Mécanique + DS \\ Niveau Brevet}
\date{}

\begin{document}
\maketitle

%====================================================
\section*{PARTIE 1 – EXERCICES}

%====================
\subsection*{★ Niveau 1}

\begin{enumerate}
\item Une voiture parcourt 100 km en 2 h. Calculer sa vitesse.
\item Donner la formule de la vitesse.
\item Un objet est immobile sur une table. Est-il en mouvement ?
\item Donner un exemple de mouvement rectiligne.
\item Donner un exemple de mouvement circulaire.
\end{enumerate}

%====================
\subsection*{★★ Niveau 2}

\begin{enumerate}[resume]
\item Un cycliste parcourt 30 km en 1h30. Calculer sa vitesse.
\item Convertir 72 km/h en m/s.
\item Un objet de 1 kg est posé sur la Terre. Calculer son poids.
\item Donner la direction et le sens du poids.
\item Citer deux forces qui peuvent s’exercer sur un objet.
\end{enumerate}

%====================
\subsection*{★★★ Niveau 3}

\begin{enumerate}[resume]
\item Une voiture parcourt 150 m en 10 s. Calculer sa vitesse.
\item Un objet de 3 kg est lâché. Calculer son poids.
\item Décrire les forces sur un livre posé.
\item Un mouvement est-il uniforme si la vitesse change ?
\item Donner un exemple de mouvement accéléré.
\end{enumerate}

%====================
\subsection*{★★★★ Niveau 4}

\begin{enumerate}[resume]
\item Une voiture roule à 90 km/h. Quelle distance en 10 s ?
\item Un objet de 5 kg → calculer le poids.
\item Un objet tombe : quelles forces ?
\item Expliquer pourquoi une voiture ralentit en freinant.
\item Donner un exemple où les forces se compensent.
\end{enumerate}

%====================
\subsection*{★★★★★ Niveau 5}

\begin{enumerate}[resume]
\item Une voiture passe de 0 à 20 m/s en 5 s. Décrire le mouvement.
\item Comparer deux vitesses : 36 km/h et 10 m/s.
\item Un objet est soumis à deux forces opposées. Que se passe-t-il ?
\item Pourquoi la ceinture de sécurité est-elle importante ?
\item Expliquer le lien entre force et changement de mouvement.
\end{enumerate}

%====================================================
\newpage
%====================================================
\newpage
\section*{DS N°1 – Voyage en voiture et sécurité routière (100 points)}

Une famille part en vacances en voiture.  
Le trajet total est de 360 km.

La voiture roule sur trois phases :
\begin{itemize}
\item Phase 1 : 120 km en 1h30
\item Phase 2 : 180 km en 2h
\item Phase 3 : 60 km en 1h
\end{itemize}

On donne : $1\ \text{km/h} = 0{,}28\ \text{m/s}$

\bigskip

\textbf{Exercice 1 – Analyse des vitesses (20 pts)}

\begin{enumerate}
\item Calculer la vitesse moyenne sur chaque phase.
\item Comparer les vitesses.
\item Quelle phase correspond à la vitesse la plus faible ?
\item Proposer une explication réaliste (trafic, route, etc.).
\end{enumerate}

\bigskip

\textbf{Exercice 2 – Vitesse globale (20 pts)}

\begin{enumerate}
\item Calculer la durée totale du trajet.
\item Calculer la vitesse moyenne sur l’ensemble du trajet.
\item Convertir cette vitesse en m/s.
\end{enumerate}

\bigskip

\textbf{Exercice 3 – Distance de réaction (20 pts)}

On considère qu’un conducteur met 1 seconde pour réagir.

\begin{enumerate}
\item Calculer la distance parcourue en 1 seconde à 90 km/h.
\item Que représente cette distance ?
\item Pourquoi est-elle importante pour la sécurité ?
\end{enumerate}

\bigskip

\textbf{Exercice 4 – Forces et freinage (20 pts)}

Lors d’un freinage, plusieurs forces s’exercent.

\begin{enumerate}
\item Citer les forces qui agissent sur la voiture.
\item Indiquer leur direction et leur sens.
\item Expliquer pourquoi la voiture ralentit.
\end{enumerate}

\bigskip

\textbf{Exercice 5 – Sécurité (20 pts)}

\begin{enumerate}
\item Expliquer le rôle de la ceinture de sécurité.
\item Relier cela au principe d’inertie.
\item Donner un exemple concret d’accident évité.
\end{enumerate}

%====================================================
\newpage
\section*{DS N°2 – Étude d’un saut en parachute (100 points)}

Un parachutiste saute d’un avion.

On distingue 3 phases :
\begin{itemize}
\item Phase 1 : chute libre (accélération)
\item Phase 2 : vitesse constante
\item Phase 3 : ouverture du parachute (ralentissement)
\end{itemize}

On donne :
\begin{itemize}
\item masse du parachutiste : $80\ \text{kg}$
\item $g = 9{,}8\ \text{N/kg}$
\end{itemize}

\bigskip

\textbf{Exercice 1 – Poids (20 pts)}

\begin{enumerate}
\item Calculer le poids du parachutiste.
\item Donner l’unité.
\item Indiquer direction et sens.
\end{enumerate}

\bigskip

\textbf{Exercice 2 – Phase 1 : chute libre (20 pts)}

\begin{enumerate}
\item Décrire le type de mouvement.
\item Pourquoi la vitesse augmente-t-elle ?
\item Quelle est la force principale ?
\end{enumerate}

\bigskip

\textbf{Exercice 3 – Phase 2 : vitesse constante (20 pts)}

\begin{enumerate}
\item Comment évolue la vitesse ?
\item Que peut-on dire des forces ?
\item Expliquer pourquoi elles se compensent.
\end{enumerate}

\bigskip

\textbf{Exercice 4 – Phase 3 : ouverture du parachute (20 pts)}

\begin{enumerate}
\item Que devient la vitesse ?
\item Quelle force apparaît fortement ?
\item Expliquer pourquoi la chute ralentit.
\end{enumerate}

\bigskip

\textbf{Exercice 5 – Application (20 pts)}

\begin{enumerate}
\item Pourquoi un parachute est-il indispensable ?
\item Donner un autre exemple où une force ralentit un mouvement.
\item Expliquer le rôle de l’air dans ce phénomène.



\end{enumerate}
%====================================================
\newpage
\section*{DS N°3 – Étude d’un trajet en train (100 points)}

Un train effectue un trajet entre deux villes distantes de 450 km.

Le mouvement du train se décompose en trois phases :
\begin{itemize}
\item Phase 1 : accélération pendant 10 minutes jusqu’à atteindre 180 km/h
\item Phase 2 : mouvement uniforme pendant 2 heures à 180 km/h
\item Phase 3 : ralentissement pendant 20 minutes jusqu’à l’arrêt
\end{itemize}

On donne : $1\ \text{km/h} = 0{,}28\ \text{m/s}$

\bigskip

\textbf{Exercice 1 – Conversion et vitesse (20 pts)}

\begin{enumerate}
\item Convertir 180 km/h en m/s.
\item Pourquoi utilise-t-on souvent le m/s en physique ?
\end{enumerate}

\bigskip

\textbf{Exercice 2 – Distance parcourue (20 pts)}

\begin{enumerate}
\item Calculer la distance parcourue pendant la phase 2.
\item Vérifier si cela correspond à la majeure partie du trajet.
\item Donner une interprétation réaliste.
\end{enumerate}

\bigskip

\textbf{Exercice 3 – Types de mouvement (20 pts)}

\begin{enumerate}
\item Identifier le type de mouvement pour chaque phase.
\item Justifier les réponses.
\end{enumerate}

\bigskip

\textbf{Exercice 4 – Forces en jeu (20 pts)}

\begin{enumerate}
\item Citer les forces qui s’exercent sur le train.
\item Expliquer pourquoi le train accélère.
\item Expliquer pourquoi il ralentit en fin de trajet.
\end{enumerate}

\bigskip

\textbf{Exercice 5 – Sécurité et inertie (20 pts)}

Lors d’un freinage brusque, les passagers sont projetés vers l’avant.

\begin{enumerate}
\item Expliquer ce phénomène.
\item Relier cela au principe d’inertie.
\item Donner un dispositif de sécurité utilisé dans les trains.
\end{enumerate}

%====================================================
\newpage
\section*{DS N°4 – Étude d’un lancer de ballon (100 points)}

Un joueur lance un ballon verticalement vers le haut.

On distingue trois phases :
\begin{itemize}
\item Phase 1 : montée du ballon
\item Phase 2 : arrêt momentané au sommet
\item Phase 3 : descente du ballon
\end{itemize}

On néglige les frottements de l’air dans une première approche.

\bigskip

\textbf{Exercice 1 – Forces (20 pts)}

\begin{enumerate}
\item Quelle est la force principale qui agit sur le ballon ?
\item Donner sa direction, son sens et son unité.
\item Cette force change-t-elle pendant le mouvement ?
\end{enumerate}

\bigskip

\textbf{Exercice 2 – Phase de montée (20 pts)}

\begin{enumerate}
\item Comment évolue la vitesse du ballon ?
\item Quel est le type de mouvement ?
\item Expliquer pourquoi le ballon ralentit.
\end{enumerate}

\bigskip

\textbf{Exercice 3 – Sommet (20 pts)}

\begin{enumerate}
\item Quelle est la vitesse du ballon au sommet ?
\item Le ballon est-il immobile à cet instant ? Justifier.
\end{enumerate}

\bigskip

\textbf{Exercice 4 – Descente (20 pts)}

\begin{enumerate}
\item Décrire le mouvement du ballon.
\item Comment évolue sa vitesse ?
\item Comparer avec la phase de montée.
\end{enumerate}

\bigskip

\textbf{Exercice 5 – Approche réelle (20 pts)}

En réalité, les frottements de l’air ne sont pas négligeables.

\begin{enumerate}
\item Quelle force supplémentaire apparaît ?
\item Quel est son effet sur le mouvement ?
\item Donner un exemple où les frottements sont importants.
\end{enumerate}
\end{document}